\section{Stacks and finer arches}\label{sec:stacks}

The arches used so far answer yes or no: Boolean or not, connected or not. Stacks, Grothendieck's \emph{champs}, suggest arches that compute.

\subsection{A classifying topos represents a stack}

The classifying topos $\Set[\mathbb T]$ of a geometric theory $\mathbb T$ represents, on the 2-category of Grothendieck topoi, the pseudo-functor that sends a topos $\mathcal E$ to the category of models of $\mathbb T$ in $\mathcal E$ \cite{caramello-tst}: to classify structures is to represent a stack of models. The simplest instance is the deck $\B G$ of a group, read in \S\ref{sec:grdecks}: its topos classifies the $G$-torsors, so it is the stack of torsors, and the gerbes of Giraud, classified by a non-abelian $H^2$ \cite{giraud}, continue the same line.

\subsection{Quotient stacks as decks with two presentations}

A quotient stack $[X/G]$ comes with two presentations: by the groupoid $X\times G\rightrightarrows X$, whose topos is that of $G$-equivariant sheaves on $X$, and as a moduli problem, a stack of families of some kind of object. Reading an invariant on both is a bridge in the sense of this note, and the invariants that suit stacks are numerical or algebraic rather than logical: the Picard group, cohomology, the points. Mumford's computation $\operatorname{Pic}(\mathcal M_{1,1})\cong\mathbb Z/12$ \cite{mumford} is the archetype: an invariant of a moduli problem obtained from an explicit presentation.

\subsection{An example from the fonds}

In folder 146 Grothendieck tabulates, for $n$ marked points on the projective line, the order of the tangent line at the last point, a line bundle on $M_{0,n}$ that descends to the quotient by $\mathfrak S_{n-1}$; he gives the orders for four points and only guesses for five and six. The edition's article on this table \cite{hua-tangent} reads the Picard group on a second presentation of the stack: $[M_{0,n}/\mathfrak S_{n-1}]$ is the stack of monic squarefree polynomials of degree $n-1$ modulo the affine group, and on that presentation $\operatorname{Pic}[M_{0,n}/\mathfrak S_{n-1}]$ comes out cyclic of order $(n-1)(n-2)$, generated by the tangent line, in characteristic zero and for every $n\ge3$. For even $n\ge6$ the group was known through the integral Chow ring computed by Edidin and Hu \cite{edidin-hu}. This is a bridge with a Picard arch, and it gives a statement that our searches did not find in print for odd $n$; the proof and its Lean formalisation are in \cite{hua-tangent}.

\subsection{Stacks in the homotopical sense}

\emph{À la poursuite des champs} \cite{maltsiniotis} is first of all about stacks in the homotopical sense, and its test categories are bridges of yet another kind: a category $A$ is a test category when presheaves on $A$ model homotopy types, and the conditions Grothendieck gives (asphericity conditions on $A$, the existence of a separating interval in $\widehat A$) read a homotopical property of $\widehat A$ on the category $A$. The deck $\Delta$ of \S\ref{sec:grdecks} is the first test category. Folders 108 (a list of questions), 118 (scratchwork on the homotopy theory of $(\mathrm{Cat})$) and 134-2 (the typescript of the first chapter of \emph{Pursuing Stacks}, with its tables of contents) work on these conditions; \S\ref{sec:test} takes up their open questions.

\subsection{A programme}

These examples suggest a family of finer arches for the generator: for a quotient stack $[X/G]$, or more generally a stack with an explicit groupoid presentation, the Picard group, the cohomology and the points read on the presentation and compared with the moduli description. Unlike the arches of \S\ref{sec:morearches}, they produce numbers and groups, so their outputs can be tested against tables, such as Grothendieck's in folder 146, and against the OEIS. We have not run the generator with them; this section only records the connection and one worked example.
